%=============================================================================
%  Physical Intelligence pi0.5 on the Unified Engine v1
%  Build: pdflatex main.tex   (article, amsmath, booktabs)
%=============================================================================
\documentclass[11pt]{article}

\usepackage[letterpaper,margin=1in]{geometry}
\usepackage{amsmath,amssymb}
\usepackage{booktabs}
\usepackage{float}
\usepackage{caption}
\usepackage{xcolor}
\usepackage[hidelinks]{hyperref}

% ---- macros -------------------------------------------------------------
\newcommand{\pizero}{$\pi_{0.5}$}
\newcommand{\ue}{Unified~Engine}
\newcommand{\iffour}{IF4}
\newcommand{\bfteen}{bf16}
\newcommand{\alveo}{Alveo~U50}
\newcommand{\numTBD}{{\color{red}--}}

\setlength{\textfloatsep}{6pt}
\setlength{\floatsep}{6pt}
\setlength{\intextsep}{6pt}
\setlength{\parskip}{2pt}

\title{\vspace{-2em}Physical Intelligence \pizero{} --- \ue{}~v1}
\author{Apex Compute}
\date{\vspace{-1em}}

\begin{document}
\maketitle

%=============================================================================
\section{\pizero{}}
\label{sec:model}
%=============================================================================
\pizero{} is a vision--language--action policy: camera views and an instruction
in, a chunk of future control steps out. Three stages run in sequence: a SigLIP
vision encoder, a Gemma-2B prefix that produces a K/V cache, and a Gemma-300M
action expert that denoises the action chunk against that cache.

The prefix is three quarters of the arithmetic (Table~\ref{tab:budget}) and
bounds compute. The expert is small but strictly sequential: ten Euler steps of
$x_{t+\Delta t}=x_t+\Delta t\,v_\theta(x_t,t,\mathrm{KV}_{\text{prefix}})$, so
it bounds latency and its error compounds across steps.

\begin{table}[H]
\centering\small
\caption{Per-inference budget at the model's true shape.}
\label{tab:budget}
\begin{tabular}{@{}llrr@{}}
\toprule
\textbf{Stage} & \textbf{Shape} & \textbf{Params} & \textbf{GFLOP} \\
\midrule
Vision (SigLIP ViT)  & 27\,L, $d$1152, 3$\times$256 tok & $\sim$0.41\,B & 655.9 \\
Prefix (Gemma-2B)    & 18\,L, $d$2048, 576 tok          & $\sim$2.0\,B  & 2332.0 \\
Expert (Gemma-300M)  & 18\,L, $d$1024, 10 steps         & $\sim$0.31\,B & 71.0 \\
\midrule
\textbf{Total} & & \textbf{$\sim$2.7\,B} & \textbf{3058.9} \\
\bottomrule
\end{tabular}
\end{table}

%=============================================================================
\section{Porting}
\label{sec:port}
%=============================================================================
The \ue{}~v1 is a soft accelerator IP, instanced eight times on one AMD Alveo U50
card at $366$\,MHz, with a $4$\,GB DRAM budget. Two constraints shape the port: the DRAM budget and
the engine's $64$-element alignment.

\textbf{Weights.} The \bfteen{} checkpoint does not fit, so every matmul weight
is 4-bit block-quantized ($\sim$8$\times$ smaller) and dequantized to \bfteen{}
at the MAC. \iffour{} buys footprint, not speed.

\textbf{Vision head dim.} Head dim $72$ is not $64$-aligned. The $Q/K/V$
projections write $80$ lanes per head ($72$ real, $8$ zero); each head is padded
to $128$ only inside the $QK^{\top}$ contraction and compacted back to $80$ for
the output projection. Bit-identical to the unpadded math.

\textbf{Prefix masking.} A masked image slot is carried as masked placeholder
rows, so the prefill stays at its full $832$ tokens. Rotary positions are
cumulative-sum-of-mask indices, not \texttt{arange}, so a masked slot does not
advance the phase. Dropping the slot instead is supported and cuts prefill by a
quarter, bit-identically, but is off in the build measured here.

\textbf{Denoise loop.} All ten steps compile into one program and one execute
call.

%=============================================================================
\section{Quantization}
\label{sec:quant}
%=============================================================================
\subsection{\iffour{}}
\iffour{} is the accelerator's 4-bit weight format. Each $64$-element block is
quantized as signed INT4 and as FP4 E2M1, and the lower-error encoding is kept.

The choice is carried by the sign bit of the block's \bfteen{} scale, which is
an absmax quotient and therefore never negative. No side table, no extra bytes.
Decode is a $2$:$1$ mux ahead of the MAC over the INT4 and FP4 paths that
already exist. The per-block search runs once, offline.

\textbf{Why a hybrid.} INT4 spends its 16 levels uniformly; FP4 spends them
densely near zero and sparsely at the tails, with a wide gap below its largest
code. Which one wins is set by the block's crest factor, the ratio of its peak to
its bulk. Dense blocks fill every uniform level; a single outlier stretches the
shared scale and starves the rest.

Table~\ref{tab:if4} shows the crossover. The extremes go as expected. The
Gaussian row is the informative one: INT4 and FP4 tie, yet the hybrid still
gains, because crest factor varies block to block within a single tensor. The
INT4 share rises with depth through the vision tower and is lowest in the token
embedding, the most heavy-tailed weight in the model.

\begin{table}[H]
\centering\small
\caption{Weight-domain SNR (dB): INT4, FP4, and the \iffour{} hybrid, with the
fraction of blocks the hybrid assigns to INT4. Synthetic blocks above,
\pizero{} checkpoint group means below.}
\label{tab:if4}
\begin{tabular}{@{}lcccc@{}}
\toprule
\textbf{Block source} & \textbf{INT4} & \textbf{FP4} & \textbf{\iffour{}} & \textbf{\% $\to$ INT4} \\
\midrule
Uniform (bounded)          & \textbf{23.12} & 19.36 & 23.12 & 99.9 \\
Gaussian                   & 19.35 & 19.52 & \textbf{20.11} & 49.5 \\
Gaussian $+$ one outlier   & 11.81 & 16.50 & \textbf{16.50} & 0.0 \\
\midrule
Action expert   & 19.50 & 19.47 & \textbf{20.25} & 54.6 \\
Prefix LM       & 19.05 & 19.47 & \textbf{20.00} & 45.2 \\
Vision encoder  & 18.31 & 19.23 & \textbf{19.65} & 37.4 \\
Token embedding & 15.39 & 18.53 & \textbf{18.65} & 17.6 \\
\bottomrule
\end{tabular}
\end{table}

Scope: the build measured below pins the INT4 branch. The hybrid's gain is a
weight-domain result and has not yet been carried end-to-end.

\subsection{The quantization floor}
Table~\ref{tab:quant} is the quantizer's cost in isolation: the same torch run,
inputs and noise at \iffour{} and at \bfteen{}, diffed per stage, with no
hardware in the loop.

Vision is nearly untouched. The prefix caches degrade with depth, V more than K.
The velocity field degrades over the last Euler steps as the integrator
accumulates error, and the final action lands at 28.5\,dB. The loss sits in the
low-magnitude action dimensions; the gripper holds above 50\,dB. SNR is a
diagnostic; the gate is closed-loop success (Sec.~\ref{sec:libero}).

\begin{table}[H]
\centering\small
\caption{\iffour{} vs.\ \bfteen{} per section (torch reference, sample~0).}
\label{tab:quant}
\begin{tabular}{@{}lcc@{}}
\toprule
\textbf{Model section} & \textbf{SNR (dB)} & \textbf{Range (dB)} \\
\midrule
Vision encoder (3 cameras)        & 25.39 & 24.3--26.6 \\
Prefix K-cache (18 layers)        & 16.57 & 13.6--19.9 \\
Prefix V-cache (18 layers)        & 11.40 & 8.1--17.0 \\
Denoise velocity $v_t$ (10 steps) & 31.51 & 22.2--34.5 \\
\midrule
\textbf{End-to-end action (10$\times$7)} & \textbf{28.47} & MSE $2.22{\times}10^{-4}$ \\
\bottomrule
\end{tabular}
\end{table}

\subsection{LIBERO}
\label{sec:libero}
LIBERO is a tabletop manipulation benchmark for language-conditioned policies,
scored by closed-loop task success. It provides two camera views; \pizero{}
takes three, so the third slot is zero-filled and masked.

An episode is many inferences. The simulator executes the first ten steps of
each chunk and requests the next, so one episode is a dozen or more full passes
on hardware. Table~\ref{tab:libero} compares the FPGA policy with the CUDA
reference over ten prompts at a fixed seed and fixed denoise noise, so both see
byte-identical observations.

\begin{table}[H]
\centering\small
\caption{LIBERO closed-loop outcomes, FPGA (\iffour{}) vs.\ CUDA (\bfteen{}):
10 prompts, one episode each, fixed seed and noise. \checkmark\ success,
$\times$ failure, number = episode length in sim steps. $\Delta$ = FPGA $-$ CUDA
steps on successes; failures run to the cap.}
\label{tab:libero}
\begin{tabular}{@{}ccccc@{}}
\toprule
\textbf{Prompt} & \textbf{FPGA (\iffour{})} & \textbf{CUDA (\bfteen{})} & \textbf{$\Delta$ steps} & \textbf{Agree} \\
\midrule
0 & \checkmark~94  & \checkmark~95  & $-1$  & \checkmark \\
1 & \checkmark~160 & \checkmark~126 & $+34$ & \checkmark \\
2 & \checkmark~123 & \checkmark~115 & $+8$  & \checkmark \\
3 & \checkmark~108 & \checkmark~109 & $-1$  & \checkmark \\
4 & $\times$~230 & $\times$~230 & --- & \checkmark \\
5$^\dagger$ & $\times$~290 & $\times$~290 & --- & \checkmark \\
6 & \checkmark~121 & \checkmark~120 & $+1$  & \checkmark \\
7 & $\times$~230 & $\times$~230 & --- & \checkmark \\
8 & \checkmark~107 & \checkmark~106 & $+1$  & \checkmark \\
9 & $\times$~230 & $\times$~230 & --- & \checkmark \\
\midrule
\textbf{Rate} & \textbf{6/10} & \textbf{6/10} & & \textbf{10/10} \\
\bottomrule
\end{tabular}
\end{table}

\noindent{\footnotesize $^\dagger$Prompt~5 is substituted from
\texttt{libero\_object}; the original \texttt{libero\_spatial} task~5 was
excluded during bring-up.}

\vspace{2pt}
\noindent The two policies agree on every prompt, and successful episodes
differ by a few steps.


%=============================================================================
\section{Throughput}
\label{sec:throughput}
%=============================================================================
All throughput results are collected on the next page
(Table~\ref{tab:summary}): NVIDIA platforms on the left, the \ue{} on the
right, with the ASIC projections in bold. All platforms run batch~1 and ten
denoise steps.

\textbf{FPGA.} Eight \ue{}~v1 instances on one \alveo{} at $366$\,MHz. Times
come from the engine's own cycle counters and are pure back-to-back execution;
programs are compiled once up front and that one-time cost is excluded.

\textbf{ASIC projection.} Each section is treated as scaling linearly with
engine count $\times$ clock. The ASIC rows scale the measured
\alveo{} times by $(1600\cdot N)/(366\cdot 8)$: $26.2\times$ for $N{=}48$
engines and $35.0\times$ for $N{=}64$ at $1.6$\,GHz. Utilization is held at the
measured FPGA value.

\textbf{GPUs.} The RTX~5070 baseline is openpi's JAX/XLA path, measured by us;
stages are jitted separately with a device sync at each boundary, so their sum
slightly exceeds the fused end-to-end time. The Jetson Orin Nano was also
measured by us with openpi's PyTorch path; its 8\,GB cannot hold the \bfteen{}
checkpoint, so we staged weights from storage and report compute only. Its
prefix is longer (968 tokens) because that path encodes all three slots and pads
the instruction to its full budget. Jetson Thor numbers are NVIDIA's published
figures; we did not measure them.

\clearpage
%-----------------------------------------------------------------------------
%  Throughput tables: one page. Left = NVIDIA, right = Unified Engine.
%-----------------------------------------------------------------------------
\begin{table}[H]
\centering\small
\caption{End-to-end latency per inference, batch~1, ten denoise steps.
ASIC rows (bold) are linear projections from the measured \alveo{}.}
\label{tab:summary}
\begin{tabular}{@{}llcrrr@{}}
\toprule
\textbf{Platform} & \textbf{Config} & \textbf{Precision} & \textbf{Latency (ms)} & \textbf{inf/s} & \textbf{Price} \\
\midrule
AMD Alveo U50           & $8\times$ \ue{}, $366$\,MHz   & \iffour{} & 13\,100 & 0.076 & \\
\textbf{\ue{} ASIC} & $48\times$ \ue{}, $1.6$\,GHz & \iffour{} & \textbf{500} & \textbf{2.00} & \\
\textbf{\ue{} ASIC} & $64\times$ \ue{}, $1.6$\,GHz & \iffour{} & \textbf{375} & \textbf{2.67} & \\
\midrule
NVIDIA Jetson Thor T5000 & MAXN, $130$\,W                & \bfteen{} & 132    & 7.58  & \$3\,499$^{*}$ \\
NVIDIA Jetson Orin Nano 8\,GB & $7$\,W mode$^{\ddagger}$ & \bfteen{} & 4\,266 & 0.234 & \$249$^{*}$ \\
NVIDIA RTX~5070          & $250$\,W                      & \bfteen{} & 130    & 7.68  & \$549 \\
\bottomrule
\end{tabular}

\vspace{2pt}
{\footnotesize $^{\ddagger}$Compute only; weight staging from storage excluded.
$^{*}$Developer-kit price. Prices are NVIDIA launch MSRP.}
\end{table}

\begin{table}[H]
\footnotesize
\setlength{\tabcolsep}{4pt}
\begin{minipage}[t]{0.485\textwidth}
\centering
\textbf{NVIDIA}\par\vspace{4pt}

\captionof{table}{RTX~5070 (\bfteen{}, openpi JAX/XLA). Per-stage times are
separately jitted; end-to-end is the fused \texttt{policy.infer}.}
\label{tab:timing-gpu}
\begin{tabular}{@{}lcc@{}}
\toprule
\textbf{Section} & \textbf{Time (s)} & \textbf{GFLOP/s} \\
\midrule
Vision encoder          & 0.019  & 33\,793 \\
Prefix (Gemma-2B)       & 0.079  & 29\,411 \\
Action expert           & 0.037  & 1\,935  \\
\midrule
\textbf{End-to-end}     & \textbf{0.130}  & \textbf{23\,502} \\
\textbf{inf/s}          & \multicolumn{2}{c}{7.68} \\
\bottomrule
\end{tabular}

\vspace{10pt}
\captionof{table}{Jetson Orin Nano 8\,GB (\bfteen{}, openpi PyTorch), $7$\,W
mode, 968-token prefix. Compute only.}
\label{tab:timing-jetson}
\begin{tabular}{@{}lcc@{}}
\toprule
\textbf{Section} & \textbf{Time (s)} & \textbf{GFLOP/s} \\
\midrule
Vision encoder          & 0.480 & 1\,366 \\
Prefix (Gemma-2B)       & 2.358 & 1\,686 \\
Action expert           & 1.428 & 53.7   \\
\midrule
\textbf{Compute}        & \textbf{4.266} & \textbf{1\,104} \\
\textbf{inf/s}          & \multicolumn{2}{c}{0.234} \\
\bottomrule
\end{tabular}

\vspace{10pt}
\captionof{table}{Jetson Thor T5000, MAXN $130$\,W. NVIDIA-published
figures, not measured by us.}
\label{tab:timing-thor}
\begin{tabular}{@{}lccc@{}}
\toprule
\textbf{Backend} & \textbf{Total (ms)} & \textbf{Model (ms)} & \textbf{Speedup} \\
\midrule
PyTorch \bfteen{} & $\sim$132 & $\sim$128 & $1.0\times$ \\
\bottomrule
\end{tabular}
\end{minipage}%
\hfill
\begin{minipage}[t]{0.485\textwidth}
\centering
\textbf{\ue{}}\par\vspace{4pt}

\captionof{table}{\alveo{} (\iffour{}), $8\times$ \ue{}~v1 at $366$\,MHz,
832-token prefix. Measured. Peak $8\times64$ MACs/cycle $= 374.8$\,GFLOP/s.}
\label{tab:timing-fpga}
\begin{tabular}{@{}lccc@{}}
\toprule
\textbf{Section} & \textbf{Time (s)} & \textbf{GFLOP/s} & \textbf{\% pk} \\
\midrule
Vision encoder          & 2.32  & 297.6 & 79 \\
Prefix (Gemma-2B)       & 10.03 & 340.0 & 91 \\
Action expert           & 0.75  & 101.9 & 27 \\
\midrule
\textbf{End-to-end}     & \textbf{13.1} & \textbf{318.0} & \textbf{85} \\
\textbf{inf/s}          & \multicolumn{3}{c}{0.076} \\
\bottomrule
\end{tabular}

\vspace{10pt}
\captionof{table}{\ue{} ASIC projection (\iffour{}), $1.6$\,GHz. Linear
scaling of Table~\ref{tab:timing-fpga} by engines$\times$clock.}
\label{tab:timing-asic}
\begin{tabular}{@{}lcc@{}}
\toprule
\textbf{Section} & \textbf{48 eng (ms)} & \textbf{64 eng (ms)} \\
\midrule
Vision encoder          & \textbf{88.4}  & \textbf{66.3}  \\
Prefix (Gemma-2B)       & \textbf{382.6} & \textbf{286.9} \\
Action expert           & \textbf{28.5}  & \textbf{21.4}  \\
\midrule
\textbf{End-to-end}     & \textbf{499.5} & \textbf{374.6} \\
\textbf{inf/s}          & \textbf{2.00}  & \textbf{2.67}  \\
Peak GFLOP/s            & 9\,830  & 13\,107 \\
Scale vs.\ U50          & $26.2\times$ & $35.0\times$ \\
\bottomrule
\end{tabular}
\end{minipage}
\end{table}
\clearpage


\end{document}
